\documentclass[10pt, letterpaper]{article}

% Packages:
\usepackage[
    ignoreheadfoot,
    top=2 cm,
    bottom=2 cm,
    left=2 cm,
    right=2 cm,
    footskip=1.0 cm,
]{geometry}
\usepackage{titlesec}
\usepackage{tabularx}
\usepackage{array}
\usepackage[dvipsnames]{xcolor}
\definecolor{primaryColor}{RGB}{0, 0, 0}
\usepackage{enumitem}
\usepackage{amsmath}
\usepackage[
    pdftitle={Junzhe Shao -- CV},
    pdfauthor={Junzhe Shao},
    pdfcreator={LaTeX with RenderCV},
    colorlinks=true,
    urlcolor=cyan
]{hyperref}
\usepackage{calc}
\usepackage{bookmark}
\usepackage{changepage}
\usepackage{iftex}

% Ensure machine-readable PDF (when using pdfTeX)
\ifPDFTeX
    \input{glyphtounicode}
    \pdfgentounicode=1
    \usepackage[T1]{fontenc}
    \usepackage[utf8]{inputenc}
\fi

\usepackage{charter}

% Settings
\raggedright
\pagestyle{empty}
\setcounter{secnumdepth}{0}
\setlength{\parindent}{0pt}
\setlength{\topskip}{0pt}
\setlength{\columnsep}{0.15cm}
\pagenumbering{gobble}
\titleformat{\section}{\bfseries\large}{}{0pt}{}[\vspace{1pt}\titlerule]
\titlespacing{\section}{-1pt}{0.3 cm}{0.2 cm}

\renewcommand\labelitemi{$\vcenter{\hbox{\small$\bullet$}}$}
\newenvironment{highlights}{
    \begin{itemize}[
        topsep=0.10 cm,
        parsep=0.10 cm,
        partopsep=0pt,
        itemsep=0pt,
        leftmargin=0 cm + 10pt
    ]
}{\end{itemize}}

\newenvironment{onecolentry}{
    \begin{adjustwidth}{0 cm + 0.00001 cm}{0 cm + 0.00001 cm}
}{\end{adjustwidth}}

\newenvironment{twocolentry}[2][]{
    \onecolentry
    \def\secondColumn{#2}
    \noindent\begin{minipage}[t]{\dimexpr\linewidth-4.7cm\relax}
}{
    \end{minipage}\hfill
    \begin{minipage}[t]{4.5cm}
        \raggedleft \secondColumn
    \end{minipage}
    \endonecolentry
}

\newenvironment{header}{
    \setlength{\topsep}{0pt}\par\kern\topsep\centering\linespread{1.5}
}{\par\kern\topsep}

% Keep original href behavior
\let\hrefWithoutArrow\href

\begin{document}
    % Simple pipe for inline separators
    \newcommand{\AND}{\unskip\cleaders\copy\ANDbox\hskip\wd\ANDbox\ignorespaces}
    \newsavebox\ANDbox
    \sbox\ANDbox{$|$}

    %====================
    % Header
    %====================
    \begin{header}
        \fontsize{25 pt}{25 pt}\selectfont Junzhe Shao

        \vspace{5 pt}

        \normalsize
        \mbox{Berkeley, CA}
        \kern 5.0 pt\AND\kern 5.0 pt
        \mbox{(+1) 858-766-1048}
        \kern 5.0 pt\AND\kern 5.0 pt
        \mbox{\hrefWithoutArrow{mailto:junzhe_shao@berkeley.edu}{junzhe\_shao@berkeley.edu}}
        \kern 5.0 pt
        \AND\kern 5.0 pt
        \mbox{\hrefWithoutArrow{https://junzheshao98.github.io/homepage/}{Website}}
    \end{header}

    \vspace{5 pt - 0.3 cm}

    %====================
    % Education
    %====================
\section{Education}

    \begin{twocolentry}{Sep 2023 --  2026 (Expected)}
        \textbf{University of California, Berkeley}\\
        Ph.D. Candidate in Biostatistics
    \end{twocolentry}

    \vspace{0.10 cm}
    \begin{twocolentry}{Sep 2021 -- May 2023}
        \textbf{Columbia University}\\
        M.S. in Biostatistics
    \end{twocolentry}

    \vspace{0.10 cm}
    \begin{twocolentry}{Sep 2016 -- Jun 2021}
        \textbf{Peking University}\\
        B.S. in Biological Science; Minor in Physics
    \end{twocolentry}

    \vspace{0.10 cm}
    \begin{twocolentry}{Sep 2019 -- Sep 2020}
        \textbf{University of California, San Diego}\\
        Visiting Student
    \end{twocolentry}

\section{Work Experience}

\begin{twocolentry}{May 2026 -- Present}
\textbf{Internship, Data Science Research}\\
ByteDance (TikTok), San Jose, CA
\end{twocolentry}

\vspace{0.1cm}
\begin{onecolentry}
    \begin{highlights}
        \item Developing cluster-based traffic splitting for recommendation experiments to improve assignment purity and reduce cross-cluster interference.
        \item Designed a Leiden loss approximating the mean squared error (MSE) objective for cluster randomization, improving the purity-precision Pareto frontier on billion-user-scale data.
        \item Developing covariate-adjustment methods for cluster-randomized experiments, drafting technical documentation, and supporting partner teams with methodology design.
    \end{highlights}
\end{onecolentry}

\vspace{0.20 cm}

\begin{twocolentry}{Jun 2025 -- May 2026}
\textbf{Internship, Product Development Data Science}\\
Genentech, South San Francisco, CA
\end{twocolentry}
    
\vspace{0.1cm}
\begin{onecolentry}
    \begin{highlights}
        \item Developed a bias-corrected Augmented Inverse Probability Weighting (AIPW) estimator for longitudinal external-control studies, leveraging double machine learning and semiparametric efficiency theory to improve causal effect estimation.
        \item Built reproducible R simulation pipelines to benchmark performance and drafted the  methods manuscript.
        \item Applied and validated on real clinical-trial data, delivering study-ready analyses.
    \end{highlights}
\end{onecolentry}

\vspace{0.10 cm}

%====================
\section{Research Experience}

\begin{twocolentry}{May 2025-Present}
    \textbf{Bias Corrected AIPW Estimator for External Control in Longitudinal RCTs}\\
\end{twocolentry}
\vspace{0.05 cm}
\begin{onecolentry}
    \begin{highlights}
        \item Developed a bias-corrected augmented inverse probability weighting (AIPW) estimator that leverages external control arms with longitudinal data.  
        \item Established bias correction at a rate governed by the number of time points, ensuring valid type I error control compared to naive direct pooling approaches. Provided non-asymptotic guarantees.
        \item Demonstrated higher statistical power compared to RCT-only estimators and other covariates adjustment methods.
    \end{highlights}
\end{onecolentry}
\vspace{0.10 cm}
\begin{twocolentry}{May 2025-Jan 2026}
    \textbf{Nonstationary Adaptive A/B Testing with State-space Models}\\
\end{twocolentry}
\vspace{0.05 cm}
\begin{onecolentry}
    \begin{highlights}
        \item Proposed a state-space modeling framework for adaptive AB testing in nonstationary environments with evolving treatment effects.
        \item Developed a Neyman-style optimal allocation policy via Kalman filtering and established asymptotic variance bounds for  ATE estimator and provide anytime valid confidence interval.
    \end{highlights}
\end{onecolentry}
\vspace{0.10 cm}
   \begin{twocolentry}{Jun 2023 -- Dec 2025}
        \textbf{Regularized Frameworks for  Adaptive Enrichment Designs}\\
    \end{twocolentry}
    \vspace{0.05 cm}
    \begin{onecolentry}
        \begin{highlights}
            \item Proposed a unified adaptive enrichment framework to maximize power for testing average treatment effects.
            \item Cast the problem as constrained optimization to unify different objectives.
            \item Calibrated selection bias and provided valid inference throughout the process.
        \end{highlights}
    \end{onecolentry}

\vspace{0.10 cm}
\begin{twocolentry}{Dec 2024 -- April 2025}
    \textbf{Enhanced Large Language Model for Predicting HIV Care Disengagement }\\
\end{twocolentry}
\vspace{0.05 cm}
\begin{onecolentry}
    \begin{highlights}
        \item Fine-tuned LLaMA 3.1 with over 4.8 million EMR records from Tanzania to predict disengagement from HIV care, ART non-adherence, and adverse outcomes.  
        \item Demonstrated superior predictive performance compared to supervised ML models and zero-shot LLMs, with strong internal and external validation results.  
        \item Provided clinically interpretable insights; expert evaluation confirmed the model's reasoning was clinically relevant in over 90\% of aligned cases.  
    \end{highlights}
\end{onecolentry}
\vspace{0.10 cm}

    \begin{twocolentry}{Nov 2021 -- Present}
        \textbf{Causal Inference for Non-Stationary Time Series Data}\\
    \end{twocolentry}
    \vspace{0.05 cm}
    \begin{onecolentry}
        \begin{highlights}
            \item Proposed a parametric generalization of synthetic control with time-varying, non-permutation-invariant weights.
            \item Applied the approach to study the impact of a mandatory certificate on COVID-19 vaccine compliance.
            \item Developing the R package \href{https://junzheshao5959.github.io/SSMimpute/}{\texttt{SSMimpute}} for state-space-model multiple imputation in non-stationary multivariate time series.
        \end{highlights}
    \end{onecolentry}

    \vspace{0.10 cm}
    \begin{twocolentry}{Jun 2022 -- Present}
        \textbf{Adaptive Design for Randomized Trials under High-dimensional Setting}\\
    \end{twocolentry}
    \vspace{0.05 cm}
    \begin{onecolentry}
        \begin{highlights}
            \item Investigated theoretical properties of covariate-balancing methods for randomized clinical trials under many covariates.
            \item Proved that a balancing criterion can be achieved asymptotically when the number of covariates is a vanishing proportion of the number of patients.
        \end{highlights}
    \end{onecolentry}

\vspace{0.10 cm}

\section{Publications \& Manuscripts in Preparation}
\begin{onecolentry}
    \textbf{Junzhe Shao}, Xiaoxuan Cai, Linda Valeri. ``\textit{Synthetic Control as Missing Data: A Unified State-Space Approach}'', Preparing for submission.
\end{onecolentry}
\vspace{0.20 cm}

\begin{onecolentry}
    \textbf{Junzhe Shao}, Jiawen Zhu, Herbert Pang. ``\textit{Leveraging Longitudinal Data in Designing Hybrid RCT with the Incorporation of Real-World Data}'', Submitted to \textit{Biostatistics}, 2026.
\end{onecolentry}
\vspace{0.20 cm}

\begin{onecolentry}
    \textbf{Junzhe Shao}, Waverley Wei, Jingshen Wang. ``\textit{Adaptive A/B Testing under Nonstationary Dynamics using State-Space Models}'',  \href{https://virtual.aistats.org/virtual/2026/poster/10811}{\textit{International Conference on Artificial Intelligence and Statistics (AISTATS)}}, 2026.
\end{onecolentry}
\vspace{0.20 cm}
\begin{onecolentry}
    \textbf{Junzhe Shao}, Aibo Gong, Juan Shen, Waverley Wei. ``\textit{A Unified Adaptive Enrichment Design for Power Enhancement}'', Major Revision at \textit{Statistica Sinica}, 2025.
\end{onecolentry}
\vspace{0.20 cm}

\begin{onecolentry}
    Waverly Wei$^\dagger$, \textbf{Junzhe Shao}$^\dagger$, Rita Qiuran Lyu, Rebecca Hemono, Xinwei Ma et al. ``\textit{Enhanced language models for predicting and understanding HIV care disengagement: a case study in Tanzania}'', \href{https://www.nature.com/articles/s41746-026-02349-3}{\textit{npj Digital Medicine}}, 2026.
\end{onecolentry}
\vspace{0.20 cm}
\begin{onecolentry}
\textbf{Junzhe Shao}, Mingzhang Yin, Xiaoxuan Cai, Linda Valeri. ``\textit{Generalized Synthetic Control Method with State-space Model}'', short version accepted at \href{https://openreview.net/pdf?id=OwyiIBIFCrn}{NeurIPS 2022 Workshop for Causality for realworld impact}. Oral Presentation at American Causal Inference Conference (ACIC) 2023, Manuscript in preparation.
\end{onecolentry}

\vspace{0.20 cm}
\begin{onecolentry}
Alton Barbehenn, Lei Shi, \textbf{Junzhe Shao}, et al. ``\textit{Rapid Biphasic Decay of Intact and Defective HIV DNA Reservoir During Acute Treated HIV Disease}'', \href{https://www.nature.com/articles/s41467-024-54116-1}{\textit{Nature Communications}}, 2024.
\end{onecolentry}
\vspace{0.20 cm}

\begin{onecolentry}
Mona Alotaibi$^\dagger$, \textbf{Junzhe Shao}$^\dagger$, Michael W. Pauciulo, et al. ``\textit{Metabolomic Profiles Differentiate Scleroderma-PAH from Idiopathic PAH and Correspond with worsened Functional Capacity}'', \href{https://www.sciencedirect.com/science/article/abs/pii/S0012369222037060}{\textit{Chest}}, 2022.
\end{onecolentry}

\vspace{0.30 cm}

\begin{onecolentry}\footnotesize
$^\dagger$\,Co-first authors.
\end{onecolentry}

\section{Skills Summary}
    \begin{onecolentry}
        \textbf{Programming:} R, Python, SQL\\
        \textbf{Software/Tools:} Spark, Ray, Hadoop, Hive, \LaTeX, Git, Bash, Slurm\\
        \textbf{Languages:} Mandarin (Native), English (Proficient)
    \end{onecolentry}

\end{document}
